% Source-derived chapter 01: Introduction
% Original source: riemann-roch-polynomials-hyperkahler-fourfolds
\chapter{Introduction}
\label{riemann-roch-polynomials-hyperkahler-fourfolds-chapter-01}
\source{riemann-roch-polynomials-hyperkahler-fourfolds}

\begin{remark}[Source section]
\label{riemann-roch-polynomials-hyperkahler-fourfolds-chapter-01-source}
\source{riemann-roch-polynomials-hyperkahler-fourfolds}
\source{riemann-roch-polynomials-hyperkahler-fourfolds}
This chapter reproduces the corresponding section of the retrieved
LaTeX source, including its definitions, statements, examples, and
proofs. It has not yet been translated into Lean.
\notready
\end{remark}

% The source section title was: Introduction
\label{sec:intro}
\source{riemann-roch-polynomials-hyperkahler-fourfolds}

We work over the complex numbers.\
A \hk\ manifold is a simply connected smooth compact K\"ahler manifold which carries a nowhere degenerate holomorphic 2-form (called a symplectic form) unique up to multiplication by a nonzero  constant.\
This important class of manifolds is a generalization in all even dimensions of K3 surfaces.

There are two key conjectures about \hk\ manifolds.

\begin{conj}
\source{riemann-roch-polynomials-hyperkahler-fourfolds}
\notready\label{conj:defoSYZ}
\source{riemann-roch-polynomials-hyperkahler-fourfolds}
Any \hk\ manifold can be deformed into a \hk\ manifold with a Lagrangian fibration.
\end{conj}

Let $2n$ be the dimension of $X$ and assume that there exists a nonzero class $\lll\in H^2(X,\Z)$ such that $\int_X\lll^{2n}=0$.\
One can deform $X$ so that $\lll$ becomes of type~$(1,1)$, hence is the first Chern class of some nontrivial holomorphic line bundle $L$ on $X$.\
For a very general deformation of this kind, $X$ has Picard number one.\
In this case, its K\"ahler cone coincides with its positive cone (\cite{HuyHK}, \cite[Proposition 3.2]{huyKahler}) and, therefore, $L$ or its dual is nef.\
Conjecture~\ref{conj:defoSYZ} is then implied by the following.

\begin{conj}[\hk\ SYZ conjecture]
\source{riemann-roch-polynomials-hyperkahler-fourfolds}
\notready\label{conj:SYZ}
\source{riemann-roch-polynomials-hyperkahler-fourfolds}
Let $X$ be a \hk\ manifold of dimension~$2n$.\
Any nontrivial nef line bundle $L$ on $X$ such that $\int_X c_1(L)^{2n}=0$ is semi-ample.
\end{conj}

The existence of a class $\lll$ as above is equivalent to the existence of a nonzero class in~$H^2(X,\Q)$ that is isotropic with respect to the Beauville--Bogomolov form $q_X$ (see equation~\eqref{fuj}).\
By Meyer's theorem, such isotropic classes exist when \mbox{$b_2(X)\geq 5$.}\

A line bundle $L$ as in Conjecture~\ref{conj:SYZ} would then define a fibration $f\colon X\to B$, with~$B$ normal.\
By results of Matsushita, $B$ is projective, $f$ is a Lagrangian hence equidimensional fibration, and any smooth fiber $X_b\coloneqq f^{-1}(b)$ is an abelian variety, endowed  with a canonical polarization which is the positive generator of the saturation of the image of the rank-1 restriction map $H^2(X,\Z)\to H^2(X_b,\Z)$ (see Section~\ref{sec:LagrangianFibration}).\ 

It is conjectured that the base $B$ of any Lagrangian fibration is  smooth, and in fact $\P^n$; note that $B$ is smooth if and only if $f$ is flat.\ 
When~$X$ is projective and $B$ is smooth, it was proved by Hwang in \cite{hw} that $B$ is $\P^n$; this was extended to the non-projective case by Greb and Lehn in \cite{gl}.\ When~$X$ is projective and~$n=2$, results of Ou in \cite{ou} and  Huybrechts and Xu  in \cite{hx} imply that $B$ is always $\P^2$.\ 

Both conjectures are true  for all known examples of \hk\ manifolds: for Conjecture~\ref{conj:defoSYZ}, see Section~\ref{subsec:SYZKnownExs}; the following stronger form of  Conjecture~\ref{conj:SYZ} follows from work of Matsushita and others.

\begin{theo}\label{thm:SYZKnownExs}
\source{riemann-roch-polynomials-hyperkahler-fourfolds}
Let $X$ be a \hk\ manifold of dimension $2n$ of either $\mathrm{K3}^{[n]}$,  generalized Kummer,  $\mathrm{OG10}$, or $\mathrm{OG6}$ deformation type.\
Let $L$ be a   nef line bundle on $X$ such that $\int_X c_1(L)^{2n}=0$ and $c_1(L)$ is primitive in $H^2(X,\Z)$.\
There exists a Lagrangian fibration $f\colon X\to \P^n$ such that $f^*\cO_{\P^n}(1)\cong L$.
\end{theo}

Indeed, a ``rational version'' of Conjecture~\ref{conj:SYZ}, namely \cite[Conjecture 1.1]{mat:isotropic}, is known for a  movable line bundle $L$ with primitive isotropic first Chern class by \cite[Corollary 1.1]{mat:isotropic} in the case of $\mathrm{K3}^{[n]}$ or generalized Kummer deformation type (based on \cite{bm,marlag,yos}), by \cite[Theorem 2.2]{mo} in the case of OG10 deformation type, and by \cite[Theorem 7.2]{mora} in the case of OG6 deformation type.\
When $L$ is nef,  we can apply \cite[Claim 3.1 and Claim 3.2]{mat:isotropic} to conclude.

One of our main results is that both conjectures are true in dimension 4 under an additional topological assumption that we now explain.\ We introduce another class $\mm\in H^2(X,\Z)$ with $q_X(\lll,\mm)>0$.\ 
The number
\begin{equation}\label{defaintro}
\source{riemann-roch-polynomials-hyperkahler-fourfolds}
    a\coloneqq\frac{1}{n!}\int_X\lll^n\mm^n
\end{equation}
is then a positive integer (see Lemma~\ref{lemmaa}).\ 
For most of this article, {\em we make the assumption $a=1$, that is, $\int_X\lll^n\mm^n=n!$} (the minimal possible value).

Assume  $\mm$ is the first Chern class of  a line bundle $M$ on $X$.\
When there is a Lagrangian fibration $f\colon X\to \P^n$ and $L=f^*\cO_{\P^n}(1)$, the restriction of $M$ to any smooth fiber is  a polarization of ``degree''~$a$.\
So the condition $a=1$ means that these  polarizations are  principal.\
This set up was the starting point of this work.\
We prove in Section~\ref{sec:LagrangianFibration}  that upon replacing $\mm$ by $\mm+r\lll$, for a suitable integer~$r$,  we can assume $q_X(\lll,\mm)=1$ and $q_X(\mm)=0$.\
By mimicking the work \cite{ort} of R\'ios Ortiz, we find that the  Huybrechts--Riemann--Roch polynomial~$P_{RR,X}$ (see  Section~\ref{subsec:RRpolynomial} for the definition) then takes the very simple form
\begin{equation}\label{eqRRintro}
\source{riemann-roch-polynomials-hyperkahler-fourfolds}
\forall k\in\Z\qquad P_{RR,X}(2k)=P_{RR,X}(q_X(k\lll+\mm))=\chi(X,L^k\ot M)=\chi(\P^n,\cO_{\P^n}(k+1))
\end{equation}
(see Theorem~\ref{th21} for more properties of~$X$).

We are naturally led to asking whether this conclusion still holds without assuming the existence of the Lagrangian fibration~$f$.

\begin{conj}
\source{riemann-roch-polynomials-hyperkahler-fourfolds}
\notready\label{conjnum}
\source{riemann-roch-polynomials-hyperkahler-fourfolds}
Let $X$ be a \hk\ manifold of dimension $2n$ with  classes $\lll,\mm\in H^2(X,\Z)$ such that
\[
\int_X\lll^{2n}=0\quad\textnormal{and}\quad\int_X\lll^n\mm^n=n!.
\]
Then the Huybrechts--Riemann--Roch polynomial of $X$  satisfies
$$\forall k\in\Z\qquad P_{RR,X}(2k)=\chi(\P^n,\cO_{\P^n}(k+1))=\binom{k+1+n}{n}.$$
\end{conj}

Conjecture \ref{conjnum} is not strictly speaking implied by Theorem~\ref{th21} and the SYZ conjecture.\ Indeed, the latter predicts that if $L$ is a nef holomorphic line bundle, {\it some positive power} of~$L$ is generated by global sections and gives a Lagrangian fibration.\
The relation~\eqref{eqRRintro} proves Conjecture~\ref{conjnum}  when $\lll$ is the first Chern class of a nef holomorphic line bundle $L$ such that~$L$ itself is generated by global sections.\
The question of whether $L$ is generated by global sections, if a power induces a Lagrangian fibration, was recently studied  in~\cite{kv}, where the authors give  a sufficient condition for this to happen; unfortunately, their result does not apply in our situation.

Our other results exclusively deal with the case of dimension $4$.\
In that dimension, there are two known deformation classes of \hk\ manifolds, namely the deformations of the Hilbert square $S^{[2]}$ of  a K3 surface $S$ (called the $\mathrm{K3}^{[2]}$ deformation type) and that of the generalized Kummer variety $\K_2(A)$ of an abelian surface $A$.\

We say that a \hk\ fourfold $X$ is \emph{of $\mathrm{K3}^{[2]}$ numerical type} if, for some K3 surface~$S$, there exists an isomorphism of abelian groups $\psi\colon H^2(X,\Z)\isomto H^2(S^{[2]},\Z)$ such that, for all $\alpha\in H^2(X,\Z)$, we have $\int_X \alpha^4 = \int_{S^{[2]}} \psi(\alpha)^4$ (this is equivalent to asking that the lattices $(H^2(X,\Z),q_X) $ and $(H^2(S^{[2]},\Z),q_{S^{[2]}}) $ be isometric and the Fujiki constants be the same).\

In \cite{og1}, O'Grady  conjectured  that a \hk\ fourfold   of  K3$^{[2]}$ numerical type is of~K3$^{[2]}$ deformation type and provided strong evidence for this statement.\
One of our results implies O'Grady's conjecture under a much weaker  topological  hypothesis.

\begin{theo}\label{thm:MainThm}
\source{riemann-roch-polynomials-hyperkahler-fourfolds}
Let $X$ be a \hk\ fourfold.\
Assume there are classes $\lll$, $\mm\in H^2(X,{\Z})$ such that $\int_X \lll^4=0$ and $\int_X\lll^2\mm^2=2$.\
Then $X$ is of~$\mathrm{K3}^{[2]}$ deformation type.
\end{theo}

\begin{coro}[O'Grady's conjecture]\label{cor:OGrady}
\source{riemann-roch-polynomials-hyperkahler-fourfolds}
A \hk\ fourfold   of  $\mathrm{K3}^{[2]}$ numerical type  is of~$\mathrm{K3}^{[2]}$ deformation type.
\end{coro}

Let us explain  how we prove Theorem~\ref{thm:MainThm}.\
The first step   is to show Conjecture~\ref{conjnum} in dimension $4$ (see Theorem~\ref{th43a}). 

\begin{theo}\label{thm:RRpoly}
\source{riemann-roch-polynomials-hyperkahler-fourfolds}
Let $X$ be a \hk\ fourfold with classes $\lll,\mm\in H^2(X,\Z)$ such that $\int_X \lll^4=0$ and $\int_X\lll^2\mm^2=2$.\
The Huybrechts--Riemann--Roch polynomial of $X$ satisfies
\[
P_{RR,X}(2k)=\chi(\P^2,\cO_{\P^2}(k+1))=\binom{k+3}{2}.
\]
Furthermore, the Fujiki constant and the Hodge and Chern numbers of $X$ are those of the Hilbert square of a K3 surface.
\end{theo}

The conclusion  of Theorem \ref{thm:RRpoly} is  weaker than being of ${\mathrm K3}^{[2]}$ numerical type.\
Theorem~\ref{thm:RRpoly} is the starting point for the  second step of the proof, in which we establish
 Conjecture~\ref{conj:SYZ} in dimension $4$  under the same  assumptions on $X$ made in Theorem~\ref{thm:MainThm}.\
More precisely, our key result is the following.

\begin{theo}\label{thm:SYZ}
\source{riemann-roch-polynomials-hyperkahler-fourfolds}
Let $X$ be a \hk\ fourfold and let $L$ be  a nef line bundle on $X$.\ Set
$\lll\coloneqq c_1(L) \in H^2(X,\Z)$ and assume that there exists $\mm\in H^2(X,\Z)$ such that     $\int_X \lll^4=0$ and $\int_X\lll^2\mm^2=2$.\
There exists a Lagrangian fibration $f\colon X\to\P^2$ with $f^*{\cO_{\P^2}}(1)\isom L$.
\end{theo}
As we mentioned earlier, \hk\ manifolds  of $\mathrm{K3}^{[n]}$ deformation type satisfy the SYZ conjecture, so Theorem \ref{thm:SYZ} is weaker than Theorem~\ref{thm:MainThm}.\  In our approach, Theorem \ref{thm:SYZ} (or rather  some weaker versions of it) is a step towards
proving   Theorem~\ref{thm:MainThm}, the precise logical relationships being   as follows.\
We first consider the case  of a very general \hk\ fourfold~$X$   satisfying the   assumptions  of Theorem \ref{thm:SYZ} and for which the class $\mm$ is also of type~$(1,1)$, that is, $\mm=c_1(M)$ for some line bundle $M$ on $X$.\
In such a case, we can assume that $\mm$ also satisfies $\int_X \mm^4=0$ and is in the boundary of the positive cone of~$X$.

A study of the effective cone of $X$ then shows that there are two cases.
\begin{itemize}
    \item The first case, mostly studied in Section \ref{sec:TwoNefIsotropic}, is when $L$ and $M$ are both nef and $L\ot M$ is ample.\
We show in  Proposition~\ref{prop:CaseC1} that in this case,
\begin{itemize}
    \item[$\circ$] either, after possibly permuting $L$ and $M$, Theorem~\ref{thm:SYZ} holds:   the linear system~$|L|$ induces a Lagrangian fibration to $\P^2$.\ But this contradicts what we had proved earlier in Section~\ref{subsec:ExceptionalDivisor}: that, if $|L|$ induces a Lagrangian fibration, $L$ and $M$ cannot be both nef;
    \item[$\circ$] or any divisor in the linear system $|L\otimes M|$ is irreducible and  the image of the rational map $\phi_{L\otimes M}\colon X\dra \P^5$ is rationally connected.\ But this contradicts   \cite[Theorem 4.2]{voisinfib}   showing that this situation cannot happen, at least when $X$ is very general as above  with Picard number $2$.
\end{itemize}
So this case in fact  does not arise (Corollary~\ref{cor:NoTwoIsotropicNef}).
    \item The second case, mostly studied in Section \ref{sec:DivisorialContraction}, is when $X$ admits a divisorial contraction and $M$ is not nef.\
    In this case, we  first prove Proposition~\ref{prop:CaseC2} (a slightly weaker version of Theorem~\ref{thm:SYZ}), namely the existence of a Lagrangian fibration $f\colon X\to\P^2$ with $f^*{\cO_{\P^2}}(1)\isom L^{k_L}$, for some positive integer~$k_L$.\
    We then use this weaker result to prove that $X$ is of  K3$^{[2]}$ deformation type (see Section~\ref{sec:OGrady}).
\end{itemize}

This completes the proof of Theorem~\ref{thm:MainThm}: as explained at the beginning of Section~\ref{sec:OGrady}, by deformation, one can always assume that the classes $\lll$ and $\mm$ are of type $(1,1)$ and the triple $(X,\lll,\mm)$ satisfies the hypotheses made above.\
But, by Theorem~\ref{thm:SYZKnownExs}, this also proves  Theorem~\ref{thm:SYZ}, once we know (by Theorem~\ref{thm:MainThm}) that $X$ is of  K3$^{[2]}$ deformation type.\ Note that our results rely, in the first case described above, on the classification work of \cite[Section 4]{voisinfib}, which extends to our context O'Grady's analysis in \cite{og1}.

We end the article with Section~\ref{sec:further}, which includes boundedness results when we fix the dimension $2n$ and the integer $a$ defined in~\eqref{defaintro}, and an (incomplete) analysis  of what happens when $n=2$ and $a$ is small.
 

\subsection*{Acknowledgements}
We would like to thank Enrico Arbarello, Ciro Ciliberto, Giovanni Mongardi, Kieran O'Grady, Gianluca Pacienza, \'Angel David R\'ios Ortiz, Giulia Sacc\`a, Jieao Song, and Chenyang Xu for useful discussions, suggestions, and references.\ We also thank the referee for useful comments.


%%%%%%%%%%%%%%%%%%%%%
